\chapter{The Formal Calculus: Grammar and Reduction Rules}
\label{app:calculus}

\chapterorient{Part~III taught the calculus the way one learns a spoken
language---by use, with intuition leading and precision following. This appendix
is the grammar book. It states, in one place and with no hand-waving, the two
pieces that \emph{define} the language: a \emph{grammar} (the legal types and
terms---what you are allowed to write down) and a \emph{small-step reduction
relation} (how an expression computes, one rewrite at a time---what those terms
\emph{mean}). Everything the main text said informally about values, records,
the state monad, \texttt{iterate\_while}, and demand-driven pruning is here as a
formal rule. You can read this appendix as the precise companion to
\cref{ch:language} (values and functions), \cref{ch:monad} (the state monad),
and \cref{ch:fold} (the fold and its pruning); each rule below carries a
back-reference to the chapter that uses it. The reward for the formality is
concrete: the whole language is a handful of reduction rules, and its one
genuinely unusual rule---demand-driven pruning---is what makes ``compute the
diagnostics'' and ``skip them for speed'' the \emph{same program}.}

\section{What a formal calculus is, and why this one is small}
\label{app:calc:orient}

A \emph{calculus}, in the sense used here, is a tiny formal language given by two
ingredients. The first is a \emph{grammar}: a finite set of rules, written in
Backus--Naur form (BNF), that says exactly which strings of symbols are legal
\emph{types} and legal \emph{terms}. The grammar is purely syntactic---it tells
you what you may write, not what it does. The second is an \emph{operational
semantics}: a \emph{reduction relation} $e \to e'$, read ``the expression $e$
steps to the expression $e'$,'' that says how a term computes by rewriting,
one small step at a time, until no rule applies and we are left with a
\emph{value}---a fully reduced answer.

This is the same machinery that defines the lambda calculus and, through it, every
typed functional language; our notation is a deliberately small dialect of it.
``Small'' is a design goal, not an accident. The calculus admits exactly what
Palace's algorithms need to be written cleanly---immutable values, records,
first-class functions, a state monad for orchestration, and one iteration
combinator---and nothing more. There is no inheritance, no exception machinery,
no mutable reference cell, no general recursion beyond the one loop combinator.
The smaller the calculus, the shorter the argument that an L3 form and its L4
re-expression compute the same thing (\cref{ch:rotations}), and the less there is
for a reader---or a checking tool---to get wrong.

\begin{keyidea}
The entire computational content of the calculus is a \emph{handful of reduction
rules}: $\beta$-reduction and \code{let} for functions, projection and
destructuring for records, three monad laws and the state effects for
orchestration, one rule for operator application, one recursive rule for
\code{iterate\_while}, and---the signature move---one rule for demand-driven
pruning. Master those and you have the formal meaning of every program in the
book. The grammar says what you may write; the reduction rules say what it means.
\end{keyidea}

\Cref{fig:appA:layered} shows the shape of the whole appendix: the grammar
defines two layers (the type language and the term language), and the reduction
relation acts on the term layer, rewriting terms to values while the type
language stands guard over what is well-formed.

\begin{figure}[t]
\centering
\begin{tikzpicture}[font=\footnotesize, node distance=5mm]
  % ===== top band: the TYPE language =====
  \node[band, fill=simBgBlue, draw=simBlue, minimum width=118mm, thick] (typesband)
    at (0,2.6) {};
  \node[font=\small\bfseries, simBlue, anchor=west] at (-5.7,2.6)
    {\textsc{type language}~~$\tau$};
  \node[font=\scriptsize, simInk, anchor=east, align=right] at (5.7,2.6)
    {\code{Scalar} \;\code{Tensor[$\sigma$]} \;\code{\{l:$\tau$\}} \;$\tau_1{\to}\tau_2$
     \;\code{Sim $\tau$} \;\code{Op[$\cdot$]} \;\code{!$\tau$}};
  % ===== middle band: the TERM language =====
  \node[band, fill=simBgGray, draw=simGray, minimum width=118mm, thick] (termsband)
    at (0,1.3) {};
  \node[font=\small\bfseries, simInk, anchor=west] at (-5.7,1.3)
    {\textsc{term language}~~$e$};
  \node[font=\scriptsize, simInk, anchor=east, align=right] at (5.7,1.3)
    {$x$ \;$\lambda x.e$ \;$e_1\,e_2$ \;\code{let} \;\code{\{l:e\}} \;\code{e.l}
     \;\code{if} \;\code{op($\bar e$)} \;\code{apply} \;\code{do} \;\code{iterate\_while}};
  % types classify terms
  \draw[depedge] (typesband.south) -- node[right, font=\scriptsize, simBlue]
    {classifies (\;$\Gamma \vdash e : \tau$\;)} (termsband.north);
  % ===== reduction acts on terms, producing values =====
  \node[product, minimum width=26mm] (val) at (0,-0.9) {a \emph{value} $v$};
  \draw[flow] (termsband.south) -- node[right, font=\scriptsize]
    {reduction $e \to e'$ (\cref{app:calc:reduction})} (val);
  \node[align=left, font=\scriptsize\itshape, simGray, text width=40mm]
    at (4.9,-0.4) {a term reduces, one small step at a time, until no rule
    applies---then it is a value.};
\end{tikzpicture}
\caption[The layered structure of the calculus]{The calculus has two grammatical
layers and one dynamics. The \emph{type language} $\tau$
(\cref{app:calc:types}) names the kinds of thing that exist---scalars, tensors,
records, functions, the simulator monad, operators. The \emph{term language} $e$
(\cref{app:calc:terms}) names the expressions you may write. A typing judgment
$\Gamma \vdash e : \tau$ says a term is well-formed and assigns it a type
(\cref{app:calc:typing}). The \emph{reduction relation} $e \to e'$
(\cref{app:calc:reduction}) is the dynamics: it rewrites a well-typed term, one
step at a time, down to a value. This appendix gives all three.}
\label{fig:appA:layered}
\end{figure}

\section{The grammar: types}
\label{app:calc:types}

We give the grammar in two halves, types first. A \emph{type} classifies values:
it is the static description of what kind of thing an expression denotes, checked
before any reduction happens. The type language is the following BNF, with each
production a way of forming a legal type $\tau$. (Throughout, $\tau$ ranges over
types, $\sigma$ over shape expressions, $l$ over field labels.)

\begin{lstlisting}[style=l4style, language=]
tau ::= Scalar | Bool | Int | Float         -- ground scalar types
      | Tensor[sigma]                        -- tensor with symbolic shape sigma
      | { l1: tau1, ..., ln: taun }          -- record (TypeScript-flavored)
      | tau1 x ... x taun                    -- tuple
      | tau1 -> tau2                          -- function
      | Sim tau                              -- simulator monad
      | Op[tau_in -> tau_out]                -- operator carrying baked-in params
      | !tau                                 -- shareable (non-linear) annotation
\end{lstlisting}

\noindent Each production earns its place. We name what it is \emph{for}, with a
one-line example, taking them in the order written.

\begin{description}[style=nextline, leftmargin=1.4em, font=\normalfont\itshape]
  \item[Ground scalars \code{Scalar}, \code{Bool}, \code{Int}, \code{Float}.] The
  atoms: a single number, a truth value, an integer, a floating-point number.
  \code{Scalar} is the generic numeric atom (a coefficient $\alpha$ in
  \code{axpy}); \code{Bool} types a predicate's answer; \code{Int} an iteration
  counter; \code{Float} an explicit floating value. \emph{Example:}
  \lstinline[style=l4style]|tol : Float|, \lstinline[style=l4style]|converged : Bool|.
  \item[Tensors \code{Tensor[$\sigma$]}.] An immutable multidimensional array
  whose shape is the symbolic expression $\sigma$ (the shape grammar is
  \cref{app:shapes}; we treat $\sigma$ as opaque here). This is the
  workhorse---every field, every vector, every operator's argument is a tensor.
  \emph{Example:} \lstinline[style=l4style]|Tensor[N]| is a length-$N$ vector;
  \lstinline[style=l4style]|Tensor[m, n]| a matrix.
  \item[Records \code{\{$l_1$: $\tau_1$, \dots, $l_n$: $\tau_n$\}}.] A bundle of
  named fields, each with its own type---the calculus's only aggregate
  (\cref{ch:records}). Solver state, configuration, and every step's output are
  records. \emph{Example:}
  \lstinline[style=l4style]|{ x: Float, residual: Float }|.
  \item[Tuples \code{$\tau_1 \times \cdots \times \tau_n$}.] A positional bundle:
  a fixed-length, heterogeneous sequence accessed by position rather than by
  name. Used where a couple of results travel together and naming them would be
  ceremony. \emph{Example:} \lstinline[style=l4style]|(Tensor[N], Scalar)|.
  \item[Functions \code{$\tau_1 \to \tau_2$}.] The type of a value that maps a
  $\tau_1$ to a $\tau_2$. The arrow associates to the right, so
  \lstinline[style=l4style]|A -> B -> C| is \lstinline[style=l4style]|A -> (B -> C)|
  (\cref{ch:language}). \emph{Example:}
  \lstinline[style=l4style]|Tensor[N] -> Scalar| is the type of \code{nrm2}.
  \item[The simulator monad \code{Sim $\tau$}.] The type of an \emph{action} that,
  when run against the simulator's state, produces a $\tau$ while threading that
  state (\cref{ch:monad}). It marks orchestration---file I/O, multi-stage
  workflows, RNG threading---as distinct from pure per-step computation.
  \emph{Example:} \lstinline[style=l4style]|get : Sim S|.
  \item[Operators \code{Op[$\tau_{\text{in}} \to \tau_{\text{out}}$]}.] A function
  value that \emph{also carries baked-in parameters}---matrix entries, a
  factorization, basis tables (\cref{ch:operators}). It is the closure ``configure
  once, apply many times'' given a first-class spelling; the brackets group the
  in/out arrow. \emph{Example:}
  \lstinline[style=l4style]|Op[Tensor[N] -> Tensor[N]]|, an assembled linear
  operator.
  \item[The shareable annotation \code{!$\tau$}.] A tag marking a value as
  \emph{shareable} (non-linear): it may be read repeatedly without being consumed.
  It distinguishes an operator's read-only baked-in parameters from the linearly
  threaded simulator state (\cref{app:calc:ownership}). \emph{Example:}
  \lstinline[style=l4style]|!LbmTables|, the velocity/weight tables closed into an
  operator.
\end{description}

\begin{notationbox}
The shape expression $\sigma$ inside \code{Tensor[$\sigma$]} has its own grammar
and its own algebra (literal axes, named axes, congruence groups
\lstinline[style=l4style]|Tensor[(S: ...)]|, rank wildcards). That grammar is the
subject of \cref{app:shapes} and is treated as opaque here: as far as the
\emph{term} calculus is concerned, \code{Tensor[$\sigma$]} is one ground type,
and shape checking is a side condition resolved by a separate shape solver. The
two appendices compose---this one gives the term dynamics, that one gives the
shape discipline---but neither depends on the internals of the other.
\end{notationbox}

\section{The grammar: terms}
\label{app:calc:terms}

The term language is what you actually write. Each production is an expression
form $e$; the reduction rules of \cref{app:calc:reduction} act on exactly these.

\begin{lstlisting}[style=l4style, language=]
e ::= x                                      -- variable
    | c                                      -- constant / literal
    | \(x: tau). e                           -- lambda abstraction
    | e1 e2                                  -- application
    | let x = e1 in e2                       -- let-binding
    | let { l1, ..., ln } = e1 in e2         -- destructuring let
    | let { l1, ..., lk, ...r } = e1 in e2   -- destructure with rest
    | { l1: e1, ..., ln: en }                -- record literal
    | { ...e0, li: ei, ..., lj: ej }         -- record spread / update
    | e.l                                    -- field access
    | (e1, ..., en)                          -- tuple literal
    | proj_i e                               -- tuple projection
    | if e0 then e1 else e2                  -- conditional
    | op(e1, ..., en)                        -- primitive operator application
    | op-with-params { p = e ; \(x: t). e }  -- operator-VALUE introducer
    | apply A e                              -- operator application (eliminator)
    | return e                               -- monadic return (Sim)
    | e1 >>= e2                              -- monadic bind
    | do { s1; ...; sn; e }                  -- do-notation block
    | get | put e | modify e                 -- state effects (Sim)
    | while_loop init cond step              -- pure tail-recursive loop
\end{lstlisting}

\noindent We group the productions by role and explain what each is for. The
order mirrors the reduction rules to come.

\paragraph{The lambda core.} A \emph{variable} \code{x} stands for a value bound
in scope; a \emph{constant} \code{c} is a literal (\code{7}, \code{true}). A
\emph{lambda} \lstinline[style=l4style]|\(x: tau). e| is an anonymous function of
one parameter \code{x} of type \code{tau} with body \code{e}---the only way to
introduce a function value. \emph{Application} \code{e1 e2} supplies an argument
to a function by juxtaposition. \emph{Let} \lstinline[style=l4style]|let x = e1 in e2|
names an intermediate: bind \code{x} to the value of \code{e1}, then evaluate
\code{e2} with \code{x} in scope. These four are the lambda calculus, and they
do everything the main text's ``values and functions'' chapter teaches
(\cref{ch:language}). \emph{Example:}
\lstinline[style=l4style]|let sq = \(x: Float). x * x in sq 5|.

\paragraph{Records.} A \emph{record literal}
\lstinline[style=l4style]|{ l1: e1, ..., ln: en }| builds a bundle of named
fields. \emph{Field access} \code{e.l} reads one field out. The two
\emph{destructuring lets} pull several fields at once---the plain form names a
fixed set of fields, the rest form
\lstinline[style=l4style]|{ l1, ..., lk, ...r }| binds some fields and collects
the remainder into \code{r}. \emph{Record spread / update}
\lstinline[style=l4style]|{ ...e0, li: ei }| builds a new record equal to
\code{e0} but with the listed fields replaced---never a mutation, always a fresh
record (\cref{ch:records}). \emph{Example:}
\lstinline[style=l4style]|{ ...s, dist: dist_next, step: s.step + 1 }|.

\paragraph{Tuples and conditionals.} A \emph{tuple literal}
\code{(e1, ..., en)} bundles values positionally; \emph{projection}
\lstinline[style=l4style]|proj_i e| reads the $i$-th component. The
\emph{conditional} \lstinline[style=l4style]|if e0 then e1 else e2| chooses a
branch on a \code{Bool}. \emph{Example:}
\lstinline[style=l4style]|if s.residual > tol then keep else stop|.

\paragraph{Operators.} \emph{Primitive operator application}
\lstinline[style=l4style]|op(e1, ..., en)| applies a built-in primitive
(\code{axpy}, \code{dot}, \code{matvec}, \dots), each governed by its own
$\delta$-rule (\cref{app:laws}). The \emph{operator-value introducer}
\lstinline[style=l4style]|op-with-params { p = e ; \(x: t). e_body }| builds a
first-class operator value: a record of closed-over \code{!}-shareable parameters
\code{p}, paired with a body lambda over the run-time argument \code{x}, of type
\lstinline[style=l4style]|Op[t -> t_out]| (\cref{ch:operators}). Its eliminator is \emph{operator
application} \code{apply A e}, which supplies the run-time argument to an operator
value \code{A}. \emph{Example:}
\lstinline[style=l4style]|apply bgkOp streamed|.

\paragraph{The monad and state effects.} \emph{Return}
\lstinline[style=l4style]|return e| injects a pure value into a \code{Sim}
action; \emph{bind} \lstinline[style=l4style]|e1 >>= e2| sequences two actions,
feeding the first's result into the second; a \emph{do-block}
\lstinline[style=l4style]|do { s1; ...; sn; e }| is the readable sugar for a chain
of binds (\cref{ch:monad}). A \code{do}-statement \code{s} is one of
\lstinline[style=l4style]|x <- e| (bind), \lstinline[style=l4style]|let x = e|
(pure binding), or a bare \code{e} (an effect of type \code{Sim ()}). The
\emph{state effects} \code{get} (read the current state), \lstinline[style=l4style]|put e|
(replace it), and \lstinline[style=l4style]|modify e| (apply a function to it) are
the only way state is touched. \emph{Example:}
\lstinline[style=l4style]|do { s <- get; put (step s) }|.

\paragraph{The loop.} \lstinline[style=l4style]|while_loop init cond step| is the
one primitive iteration form---a pure tail-recursive loop. The book's
\code{iterate\_while} (\cref{ch:fold}) is the named, trajectory-collecting
specialization of it, and its small-step rule is \cref{app:calc:iterate} below.
This is the \emph{only} unbounded control construct in the calculus; everything
else terminates structurally.

\begin{keyidea}
The grammar is closed: there is no production for a mutable variable, no
assignment, no \code{goto}, no general recursion outside \code{while\_loop}.
Every term is built from this finite list, and every term's behavior is fixed by
the reduction rules of the next section. The expressive reach---six simulators
from one vocabulary---comes from \emph{first-class functions and operators}
(lambda, \code{apply}, \code{Op}) combined with the single loop, not from a large
language.
\end{keyidea}

\section{Ownership: three categories the syntax enforces}
\label{app:calc:ownership}

Before the dynamics, one static distinction the grammar quietly encodes, because
several reduction rules below depend on it. The calculus sorts every value into
exactly one of three \emph{ownership categories}, and the sort is syntactic---you
can read it off the form.

\begin{description}[style=nextline, leftmargin=1.4em, font=\normalfont\itshape]
  \item[Simulator state.] Values threaded through the \code{Sim} monad, touched
  only via \code{get} / \code{put} / \code{modify}. Treated as \emph{linear}:
  a \code{Sim} action consumes the current state and produces the next, so the old
  state is shadowed and aliasing is structurally impossible. There is exactly one
  state slot in scope at any point in a \code{do} block.
  \item[Operator internal parameters.] Values closed over inside an
  \code{Op[$\cdot$]} and tagged \code{!} (shareable). Read-only across a solve;
  they never evolve. Matrix entries, factorizations, basis and lattice tables,
  time-step constants.
  \item[Ephemeral intermediates.] Values bound by a \code{let} or a \code{do}
  binding. Local in scope; they participate in neither state evolution nor
  operator closure. The residual computed mid-iteration, the CG search direction,
  a predictor in a multi-step integrator.
\end{description}

\noindent Linearity is tracked only to enforce the first-versus-second
distinction---which slot owns the canonical current value. It is \emph{not}
tracked at the tensor level: every tensor is immutable, so two names for one
tensor can never surprise each other, and \code{clone()} is meaningless
(\cref{ch:language}). The \code{!} annotation is the syntactic marker; inside an
\code{Op[$\cdot$]}, all closed-over data is \code{!}-tagged by construction. We
will see this category distinction surface in the state-effect rules
(\cref{app:calc:state}) and the operator-application rule
(\cref{app:calc:apply}).

\section{Small-step operational semantics}
\label{app:calc:reduction}

We now give the dynamics: the reduction relation $e \to e'$, presented as a set
of \emph{inference rules}. An inference rule is read as ``\emph{if} the premises
above the line hold, \emph{then} the conclusion below the line holds.'' For most
of our reduction rules there are no premises---they are \emph{axioms}, an
unconditional rewrite, drawn as a conclusion with an empty line above it (or
simply stated as an equation $e \to e'$). The pruning rule of
\cref{app:calc:pruning} is the one rule with a genuine premise.

Throughout, $v$ ranges over \emph{values}---fully reduced terms that no rule
rewrites further (a literal, a lambda, a fully-built record of values, an operator
value). The notation $e[v/x]$ means ``the term $e$ with every free occurrence of
the variable $x$ replaced by $v$''---capture-avoiding substitution, the single
operation underlying nearly every rule.

\Cref{fig:appA:reduction} traces a small expression all the way down to a value
by these rules, to fix the idea of ``reduce one step at a time'' before we
catalogue the steps.

\begin{figure}[t]
\centering
\begin{tikzpicture}[font=\footnotesize, node distance=3mm]
  \node[opbox, align=left, minimum width=104mm] (e0) at (0,0)
    {\lstinline[style=l4style]|let r = axpy 2 x y in r.|... \quad\textit{(start: a let over an operator app)}};
  \node[font=\scriptsize, simBlue, right=24mm of e0.east, anchor=west] (l0) {};
  \node[opbox, align=left, minimum width=104mm, below=8mm of e0] (e1)
    {\lstinline[style=l4style]|let r = (2 .* x .+ y) in r.|... \quad\textit{($\delta$-rule for \code{axpy}: unfold the primitive)}};
  \node[opbox, align=left, minimum width=104mm, below=8mm of e1] (e2)
    {\lstinline[style=l4style]|let r = v_r in r.|... \quad\textit{(the tensor arithmetic evaluates to a value $v_r$)}};
  \node[opbox, align=left, minimum width=104mm, below=8mm of e2] (e3)
    {\lstinline[style=l4style]|v_r.|... \quad\textit{(\textbf{let} rule, \S\ref{app:calc:lambda}: substitute $v_r$ for \code{r})}};
  \node[product, align=left, minimum width=104mm, below=8mm of e3] (e4)
    {a \emph{value} --- no rule applies; reduction halts.};
  % step arrows with rule labels
  \draw[flow] (e0) -- node[right, font=\scriptsize, simBlue]{$\to$ \;$\delta$} (e1);
  \draw[flow] (e1) -- node[right, font=\scriptsize, simBlue]{$\to^{*}$ \;(arith)} (e2);
  \draw[flow] (e2) -- node[right, font=\scriptsize, simBlue]{$\to$ \;\textsf{let}} (e3);
  \draw[flow] (e3) -- node[right, font=\scriptsize, simBlue]{$\to^{*}$ \;\textsf{proj}} (e4);
\end{tikzpicture}
\caption[A term reducing step by step]{Reduction is rewriting, one small step at
a time. An expression is repeatedly matched against the reduction rules; each
match rewrites it to a simpler term ($\to$ is one step, $\to^{*}$ several), until
no rule applies and the term is a \emph{value}. Here a \code{let} over an
operator application unfolds the \code{axpy} $\delta$-rule (\cref{app:laws}),
evaluates the tensor arithmetic, then fires the \code{let} rule
(\cref{app:calc:lambda}) to substitute the result. The rule fired at each step is
named on the arrow. The order is forced only by data dependency
(\cref{app:calc:confluence}); the answer does not depend on which legal step we
take first.}
\label{fig:appA:reduction}
\end{figure}

\subsection{The lambda core: \texorpdfstring{$\beta$}{beta} and \texttt{let}}
\label{app:calc:lambda}

The two foundational rules. \emph{$\beta$-reduction} is function
application---supply an argument to a lambda by substituting it into the body.
\emph{The \code{let} rule} names a value---bind a value to a name by substituting
it into the body. They are the same operation (substitution) wearing two faces.

\[
\begin{array}{rcll}
(\lambda(x:\tau).\, e)\ v &\;\to\;& e[v/x] & \text{($\beta$)} \\[2pt]
\textsf{let}\ x = v\ \textsf{in}\ e &\;\to\;& e[v/x] & \text{(\textsf{let})}
\end{array}
\]

\noindent \emph{Intuition.} $\beta$ says: to apply
\lstinline[style=l4style]|(\(x: Float). x * x)| to \code{5}, erase the lambda and
replace \code{x} by \code{5} in the body, giving \lstinline[style=l4style]|5 * 5|.
The \code{let} rule says the same for naming: \lstinline[style=l4style]|let y = 6 in y + 1|
becomes \lstinline[style=l4style]|6 + 1|. There are no boxes and no overwrites
(\cref{ch:language}): a name is a stand-in for a value, and reduction simply
substitutes the value back in. Currying and partial application
(\cref{ch:language}) are nothing more than $\beta$ applied one argument at a time:
\lstinline[style=l4style]|add 3 4| is \lstinline[style=l4style]|(add 3) 4|, two
$\beta$-steps.

\subsection{Records: projection, destructuring, spread}
\label{app:calc:records}

Records have three rules: \emph{projection} reads one field; \emph{destructuring}
binds several fields at once; \emph{spread} builds an updated copy. All operate on
a record of \emph{values} $\{l_1\!:\!v_1, \dots, l_n\!:\!v_n\}$.

\[
\begin{array}{rcll}
\{l_1\!:\!v_1,\dots,l_n\!:\!v_n\}.l_k &\;\to\;& v_k & \text{(proj)} \\[3pt]
\textsf{let}\ \{l_1,\dots,l_n\} = \{l_1\!:\!v_1,\dots,l_n\!:\!v_n\}\ \textsf{in}\ e
  &\;\to\;& e[v_1/l_1,\dots,v_n/l_n] & \text{(destr)} \\[3pt]
\{\dots\{l_1\!:\!v_1,\dots,l_n\!:\!v_n\},\ l_k\!:\!w\}
  &\;\to\;& \{l_1\!:\!v_1,\dots,l_k\!:\!w,\dots,l_n\!:\!v_n\} & \text{(spread)}
\end{array}
\]

\noindent \emph{Intuition.} Projection \code{r.x} pulls the named field's value
straight out. Destructuring is projection done in bulk: it binds each field name
to its value in one stroke, so
\lstinline[style=l4style]|let { x, residual } = s in ...| makes both fields
available without writing \code{s.x} and \code{s.residual} everywhere. Spread is
the immutable update: \lstinline[style=l4style]|{ ...s, x: x' }| produces a
\emph{new} record identical to \code{s} except in field \code{x}---the original
\code{s} is untouched, exactly as immutability requires (\cref{ch:records}). The
``destructure with rest'' form
\lstinline[style=l4style]|let { l1, ..., lk, ...r } = ...| is the same destr rule
with \code{r} bound to the record of remaining fields; we elide its rule as a
routine variant. The spread rule is the formal content of every
``\lstinline[style=l4style]|{ ...s, dist: dist_next }|'' you have seen: it is how
a step ``updates'' state without mutating it.

\subsection{The monad: \texttt{return}, \texttt{>>=}, and \texttt{do}}
\label{app:calc:monad}

Orchestration is sequenced by the three \emph{monad laws}. These govern how
\code{Sim} actions compose; they are the precise meaning of a \code{do}-block
(\cref{ch:monad}).

\[
\begin{array}{rcll}
\textsf{return}\ v \,\mathbin{\texttt{>>=}}\, f &\;\to\;& f\ v & \text{(left identity)} \\[2pt]
m \,\mathbin{\texttt{>>=}}\, \textsf{return} &\;\to\;& m & \text{(right identity)} \\[2pt]
(m \,\mathbin{\texttt{>>=}}\, f) \,\mathbin{\texttt{>>=}}\, g
  &\;\to\;& m \,\mathbin{\texttt{>>=}}\, (\lambda x.\, f\ x \,\mathbin{\texttt{>>=}}\, g)
  & \text{(associativity)}
\end{array}
\]

\noindent \emph{Intuition.} \emph{Left identity}: injecting a value with
\code{return} and then binding it is just calling the continuation on the
value---\code{return} adds no effect. \emph{Right identity}: binding an action to
\code{return} changes nothing---returning the result is a no-op. \emph{Associativity}:
how you parenthesize a chain of binds does not matter, which is exactly why the
flat \code{do}-block notation is well-defined---\lstinline[style=l4style]|do { x <- m; y <- n; k }|
is unambiguous because the binds re-associate freely. A \code{do}-block desugars
to a right-nested chain of \code{>>=}: \lstinline[style=l4style]|do { x <- m; rest }|
is \lstinline[style=l4style]|m >>= (\x -> do { rest })|, and a pure
\lstinline[style=l4style]|let x = e; rest| desugars to
\lstinline[style=l4style]|(\x -> do { rest }) e|. The three laws guarantee these
desugarings all mean the same thing.

\subsection{State effects}
\label{app:calc:state}

The state effects \code{get} / \code{put} / \code{modify} are the only way the
threaded simulator state (\cref{app:calc:ownership}, category one) is touched.
Two rewrite rules relate them, capturing that \code{modify} is exactly
``get, transform, put'':

\[
\begin{array}{rcl}
\textsf{do}\ \{\, x \leftarrow \textsf{get};\ \textsf{put}\ f(x);\ \textsf{return}\ ()\,\}
  &\;\to\;& \textsf{modify}\ f \\[3pt]
\textsf{modify}\ (f \circ g) &\;\to\;& \textsf{do}\ \{\,\textsf{modify}\ g;\ \textsf{modify}\ f\,\}
\end{array}
\]

\noindent \emph{Intuition.} The first rule says \code{modify f} is shorthand for
reading the state, applying \code{f}, and writing it back---there is no separate
``read-modify-write'' hazard because the state slot is linear and the old value is
shadowed. The second rule says modifying by a composed function $f \circ g$ is the
same as modifying by $g$ then by $f$: state updates compose in application order.
Together they let the orchestration layer be rewritten freely between the explicit
\code{get}/\code{put} form and the compact \code{modify} form without changing
meaning. Crucially, per-algorithm state evolution does \emph{not} need these---a
solver's state threads through \code{iterate\_while} as a pure carry
(\cref{app:calc:iterate}); \code{Sim} recedes to genuine orchestration (I/O,
multi-stage workflows).

\subsection{Operator application}
\label{app:calc:apply}

An operator value pairs closed-over parameters \code{p} with a body lambda
$\lambda x.\, e$ (\cref{app:calc:terms}, the operator-value introducer). Applying
it substitutes \emph{both} the parameters and the run-time argument into the body:

\[
\textsf{apply}\ (\textsf{op-with-params}\ p,\ \lambda x.\, e)\ v
  \;\to\; e[p/\textsf{params},\ v/x]
\]

\noindent \emph{Intuition.} This is $\beta$-reduction with a twist: an ordinary
lambda has only its parameter to substitute, but an operator value also carries a
record of baked-in parameters (the matrix, the factorization, the basis tables),
and applying it substitutes those \emph{and} the run-time argument
simultaneously. The notation $e[p/\textsf{params}]$ means the body's references to
its closed-over parameters are replaced by the actual closed-over values \code{p}.
This is the formal content of ``build once with configuration baked in, apply many
times'' (\cref{ch:operators}): the \code{op-with-params} introducer is the
``build'' (it captures \code{p} once), and each \code{apply A v} is one ``apply.''
The parameters \code{p} are \code{!}-shareable
(\cref{app:calc:ownership}, category two), so the same operator value may be
applied to many arguments with no copying and no consumption.

\subsection{Primitive \texorpdfstring{$\delta$}{delta}-rules}
\label{app:calc:delta}

Each primitive operator---\code{axpy}, \code{dot}, \code{matvec}, \code{conv},
\code{reduce}, \dots---comes with its own \emph{$\delta$-rule}, an axiom giving
its action on value arguments. We do not enumerate them here: the primitive
semantics live in the concept library and the algebraic-law appendix
(\cref{app:laws}), where each $\delta$-rule can be checked against a reference
implementation. The calculus assumes only that every primitive \code{op} comes
\emph{with} such a rule, so that \lstinline[style=l4style]|op(v1, ..., vn)|
reduces to a value. \Cref{fig:appA:reduction} fired the \code{axpy} $\delta$-rule
as its first step. The point of structure here is the division of labor: the
\emph{shape} of computation (functions, records, loops, pruning) is fixed by this
appendix; the \emph{arithmetic} of each primitive is fixed by its $\delta$-rule
elsewhere.

\subsection{The \texttt{iterate\_while} rule}
\label{app:calc:iterate}

Here is the heart of the dynamics: the small-step rule for the iteration
combinator the whole book turns on (\cref{ch:fold}). Given a carry $a$, a
predicate $p : \alpha \to \code{Bool}$, and a step
$f : \alpha \to \{\code{state}\!:\!\alpha, \dots e\}$, one unrolling step of
\code{iterate\_while} is:

\[
\begin{aligned}
\textsf{iterate\_while}\ a\ p\ f \;\to\;\ &\textsf{if}\ p(a)\\
  &\quad\textsf{then}\ \textsf{let}\ \{\textsf{state}\!:\!a',\ \dots e\} = f(a)\ \textsf{in}\\
  &\qquad \textsf{let}\ \{\textsf{final\_state},\ \textsf{trajectory}\}
        = \textsf{iterate\_while}\ a'\ p\ f\ \textsf{in}\\
  &\qquad \{\textsf{final\_state},\ \textsf{trajectory}\!:\
        [\{\dots e\}] \mathbin{+\!\!+} \textsf{trajectory}\}\\
  &\quad\textsf{else}\ \{\textsf{final\_state}\!:\ a,\ \textsf{trajectory}\!:\ [\,]\,\}
\end{aligned}
\]

\noindent \emph{Intuition.} Read it as a recursion that unrolls one loop step per
firing. Test the predicate on the \emph{current} carry $a$. If it is false, stop:
the answer is $a$ as \code{final\_state} with an empty \code{trajectory}. If it is
true, run one step $f(a)$, which yields the next carry $a'$ together with this
step's per-step extras $e$; recurse on $a'$; and prepend $e$ to whatever
trajectory the recursion returns. Two structural facts are baked into the rule and
are load-bearing in \cref{ch:fold}. First, the predicate fires \emph{before} any
step (the \code{while}, not \code{do}/\code{while}, convention): an
already-converged input does zero work. Second, the predicate reads only the carry
$a$---its type is $\alpha \to \code{Bool}$---so any termination quantity must be
\emph{folded into the carry}; letting the predicate read the step's extras would
be a circular dependency, and the type forbids it. The recursion sits in tail
position, so an implementation realizes it as an ordinary loop.

\Cref{fig:appA:unroll} draws one firing of the rule as the operation it really
is: a single step splits into a \emph{carry update} (the next $a'$, threaded
forward) and a \emph{trajectory extension} (the extras $e$, prepended to the
list).

\begin{figure}[t]
\centering
\begin{tikzpicture}[font=\footnotesize, node distance=7mm]
  % the rule as: one application splits into carry-update + trajectory-extend
  \node[stage, minimum width=40mm, align=center] (lhs) at (0,0)
    {\lstinline[style=l4style]|iterate_while a p f|};
  \node[font=\scriptsize, simBlue] (arrow) at (3.3,0) {$\to$ \;(one unrolling)};
  % predicate gate
  \node[diamond, draw=simBlue, fill=simBgBlue, aspect=1.7, inner sep=1pt,
        font=\scriptsize, text=simInk] (p) at (6.7,0) {$p(a)$?};
  \draw[flow] (lhs.east) -- (p.west);
  % FALSE branch
  \node[product, minimum width=34mm, align=center] (stop) at (6.7,-2.2)
    {\lstinline[style=l4style]|{ final_state: a, trajectory: [] }|};
  \draw[backflow] (p.south) -- node[right, font=\scriptsize, simRust]{no} (stop.north);
  % TRUE branch: run the step
  \node[opbox, minimum width=20mm] (step) at (10.4,0) {\lstinline[style=l4style]|f(a)|};
  \draw[flow] (p.east) -- node[above, font=\scriptsize, simGreen]{yes} (step.west);
  % step splits into carry-update and trajectory-extend
  \node[data, fill=white, draw=simGray, minimum width=20mm] (carry) at (10.4,-1.5)
    {carry update: $a'$};
  \node[data, fill=simBgGold, draw=simGold, minimum width=26mm] (traj) at (10.4,-2.9)
    {trajectory ext.: $[\{\dots e\}]\mathbin{+\!\!+}\cdots$};
  \draw[flow] (step.south) -- node[right, font=\scriptsize]{\code{.state}} (carry.north);
  \draw[refedge, simGold!80!black] (step.south west) to[out=240,in=80] (traj.north);
  % the recursion arrow: carry update feeds the recursive call
  \draw[flow, simBlue, densely dashed] (carry.west) to[out=180,in=270]
    node[below, font=\scriptsize, simBlue, pos=0.4]{recurse on $a'$} (lhs.south);
\end{tikzpicture}
\caption[One unrolling step of \texttt{iterate\_while}]{The
\code{iterate\_while} small-step rule (\cref{app:calc:iterate}), drawn as the
single operation it performs. The predicate $p(a)$ gates the firing. On
\textcolor{simRust}{no}, the loop stops, surfacing the current carry as
\code{final\_state} with an empty trajectory. On \textcolor{simGreen}{yes}, one
step \code{f(a)} runs and \emph{splits in two}: its \code{.state} field becomes
the next carry $a'$ (threaded forward, and fed back into the recursive call,
dashed), while its extras $e$ are prepended to the growing
\code{trajectory}. This is the formal version of the carry-update-plus-append
picture of \cref{ch:fold}: every iterative algorithm in the book is this rule
fired repeatedly, with ``step'' replaced by the algorithm's per-step kernel.}
\label{fig:appA:unroll}
\end{figure}

The combinator has a sugared special case,
\lstinline[style=l4style]|iterate_while_pure init cont step|, for steps with no
extras: it is \lstinline[style=l4style]|(iterate_while init cont (\a -> { state: step a })).final_state|,
the general rule with $\code{extras} = ()$ and the trajectory always empty
(\cref{ch:fold}). Total correctness---that the loop halts---is an obligation on
the \emph{algorithm} that uses the rule (a bounded \code{max\_it} folded into the
carry, or a convergence proof), not a guarantee of the rule itself; the type
promises nothing about termination.

\subsection{Demand-driven pruning: the signature move}
\label{app:calc:pruning}

The calculus has one rule unlike anything in an ordinary operational
semantics, and it is the most important rule in the appendix. It formalizes the
demand-driven pruning of \cref{ch:fold}: an output that no consumer reads is never
computed. It is the only rule here with a real premise.

The calculus is a \emph{graph-evaluation language}: operators and aggregate
operations produce \emph{records} of outputs, and at solve time only the outputs
that some root consumer transitively reaches are computed. Observation propagates
structurally through bindings---if a body reads field $l_k$, then $l_k$ is
observed in the producing record; if the body never references $l_k$ (uses a
\code{\_} placeholder, or simply omits it), $l_k$ is unobserved. The pruning rule
is then an equivalence:

\[
\dfrac{\;\text{output } l_k \text{ of } \textsf{op} \text{ is not observed in } \textsf{body}\;}
{\;
  \textsf{let}\ \{l_1\!:\!x_1, \dots, l_k\!:\!\_, \dots, l_n\!:\!x_n\}
    = \textsf{op}(\bar{e})\ \textsf{in}\ \textsf{body}
  \;\equiv\;
  \textsf{let}\ \{l_i\!:\!x_i\}_{i \neq k} = \textsf{op}_{\neg k}(\bar{e})\ \textsf{in}\ \textsf{body}
\;}
\]

\noindent where $\textsf{op}_{\neg k}$ is the subgraph of $\textsf{op}$ that does
\emph{not} compute the $l_k$-output---formally, the dependency-closure of the
remaining outputs over $\textsf{op}$'s body. \emph{Intuition.} The premise (above
the line) is a fact about \emph{use}: nothing downstream reads field $l_k$. Given
that fact, the conclusion (below the line) licenses rewriting the producer to a
strictly smaller operator that omits $l_k$ entirely, along with every
sub-computation that fed \emph{only} into $l_k$. This is dead-code elimination,
lifted from a compiler's view of a function body to the evaluation of the whole
computation graph. The rule reads as a subgraph-derivation: from the full graph
and a non-observation premise, derive the pruned subgraph.

\Cref{fig:appA:pruning} draws the derivation as the graph transformation it is:
the full dependency graph on the left, and on the right the subgraph that survives
when one output is unread, with the dead nodes and edges erased.

\begin{figure}[t]
\centering
\begin{tikzpicture}[font=\footnotesize, node distance=6mm]
  % ============ LEFT: full graph (all outputs observed) ============
  \begin{scope}
    \node[font=\small\bfseries, simBlue] at (0,2.7) {full graph: every output observed};
    \node[data, minimum width=14mm] (in) at (0,1.8) {inputs $\bar e$};
    \node[opbox, minimum width=12mm] (core) at (-1.2,0.7) {core};
    \node[opbox, minimum width=14mm] (extra) at (1.4,0.7) {extra calc};
    \draw[flow] (in.south) -- (core.north);
    \draw[flow] (in.south) -- (extra.north);
    \node[product, minimum width=12mm] (o1) at (-1.2,-0.7) {$l_1$};
    \node[product, fill=simBgGold, draw=simGold, minimum width=12mm] (ok) at (1.4,-0.7) {$l_k$};
    \draw[flow] (core.south) -- (o1.north);
    \draw[flow] (extra.south) -- (ok.north);
    % both consumed
    \draw[refedge] (o1.south) -- ++(0,-0.55) node[below, font=\scriptsize]{read};
    \draw[refedge] (ok.south) -- ++(0,-0.55) node[below, font=\scriptsize, simGold!80!black]{read};
  \end{scope}
  \draw[simGray!50, thick] (3.4,-1.6) -- (3.4,2.9);
  % the rule arrow with the premise
  \node[align=center, font=\scriptsize, simBlue] at (4.5,1.0)
    {premise:\\$l_k$ \emph{not}\\observed\\[2pt]$\Longrightarrow$};
  % ============ RIGHT: pruned subgraph ============
  \begin{scope}[xshift=8.4cm]
    \node[font=\small\bfseries, simGreen] at (0,2.7) {pruned subgraph: \code{op}$_{\neg k}$};
    \node[data, minimum width=14mm] (in2) at (0,1.8) {inputs $\bar e$};
    \node[opbox, minimum width=12mm] (core2) at (-1.2,0.7) {core};
    % the extra calc is dead (greyed)
    \node[opbox, draw=simGray!40, fill=white, text=simGray!55, minimum width=14mm]
      (extra2) at (1.4,0.7) {\textcolor{simGray!55}{extra calc}};
    \draw[flow] (in2.south) -- (core2.north);
    \draw[flow, simGray!35] (in2.south) -- (extra2.north);
    \node[product, minimum width=12mm] (o12) at (-1.2,-0.7) {$l_1$};
    \node[product, draw=simGray!40, fill=white, text=simGray!55, minimum width=12mm]
      (ok2) at (1.4,-0.7) {\textcolor{simGray!55}{$l_k$}};
    \draw[flow] (core2.south) -- (o12.north);
    \draw[flow, simGray!35] (extra2.south) -- (ok2.north);
    \draw[refedge] (o12.south) -- ++(0,-0.55) node[below, font=\scriptsize]{read};
    \node[font=\scriptsize, simGray!60] at (1.4,-1.45) {pruned};
  \end{scope}
\end{tikzpicture}
\caption[Demand-driven pruning as graph dead-code elimination]{The pruning rule
(\cref{app:calc:pruning}) as a subgraph derivation. \emph{Left}: the full
operator graph---two outputs $l_1$ and $l_k$, each produced by its own
sub-computation, both read by a consumer. \emph{Right}: the premise ``$l_k$ is
not observed'' fires the rule, deriving $\textsf{op}_{\neg k}$---the subgraph that
omits $l_k$ \emph{and} the \emph{extra calc} node that fed only into it (greyed,
erased). The surviving graph computes exactly the observed outputs. This is the
calculus's signature move: no verbosity flag, no conditional branch, no Writer
monad---one operator definition, and what materializes is decided by what the
consumer reads. Applied inside \code{iterate\_while}, it is why an unread
\code{trajectory} costs nothing (\cref{ch:fold}).}
\label{fig:appA:pruning}
\end{figure}

This single rule subsumes three mechanisms other systems need separately: a
\emph{Writer monad} (logging/metric accumulation), a \emph{verbosity flag}
(``if verbose, also compute the residual''), and a \emph{conditional branch}
around each optional output. In the calculus none of those exist as separate
concepts: every operator unconditionally declares all its outputs as record
fields, every consumer reads what it needs, and pruning removes the rest. An
algorithm is written \emph{once} and is correct whether or not its optional
outputs are consumed---which is why no solver chapter in this book presents a
``debug version'' beside a ``fast version'' (\cref{ch:fold}).

\begin{keyidea}
Demand-driven pruning is the calculus's one genuinely unusual rule, and its whole
payoff. ``Compute the residual history for monitoring'' and ``skip it for speed''
are not two algorithms but \emph{one}, distinguished only by what a downstream
consumer reads. The rule is graph dead-code elimination given as a formal
equivalence: a non-observation premise licenses replacing a producer by the
subgraph of its observed outputs. It is what lets the framework write each solver
once and use it for both diagnostics and production.
\end{keyidea}

\section{Typing in one paragraph}
\label{app:calc:typing}

The grammar and reduction rules are the load-bearing content; the type system is
standard typed-$\lambda$ and we state only its shape. A typing judgment
$\Gamma \vdash e : \tau$ reads ``in context $\Gamma$ (a set of variable-to-type
assumptions), the term $e$ has type $\tau$.'' The selected rules are the usual
ones---a variable has the type the context assigns it; an application
$e_1\,e_2$ requires $e_1 : \tau_1 \to \tau_2$ and $e_2 : \tau_1$ and yields
$\tau_2$; a record literal is typed field-by-field; a spread preserves the record
type; \code{return} lifts a $\tau$ to \code{Sim $\tau$}; and \code{>>=} sequences
\code{Sim $\alpha$} with $\alpha \to \code{Sim $\beta$}$ to give \code{Sim $\beta$}:

\begin{gather*}
\dfrac{}{\Gamma, x{:}\tau \vdash x : \tau}\;(\textsf{var})
\qquad
\dfrac{\Gamma \vdash e_1 : \tau_1 \to \tau_2 \quad \Gamma \vdash e_2 : \tau_1}
      {\Gamma \vdash e_1\,e_2 : \tau_2}\;(\textsf{app})
\\[6pt]
\dfrac{\Gamma \vdash m : \textsf{Sim}\,\alpha \quad \Gamma, x{:}\alpha \vdash k : \textsf{Sim}\,\beta}
      {\Gamma \vdash m \mathbin{\texttt{>>=}} \lambda x.k : \textsf{Sim}\,\beta}\;(\textsf{bind})
\end{gather*}

\noindent The one place the type system does real, distinctive work is
\emph{shapes}: a primitive's signature carries symbolic tensor shapes
(\lstinline[style=l4style]|axpy : Scalar -> Tensor[s] -> Tensor[s] -> Tensor[s]|),
and the requirement that two shapes ``match'' is a side condition discharged by a
separate shape solver, not by these rules. That discipline---congruence groups,
rank wildcards, the named-axis algebra---is the entire subject of
\cref{app:shapes}. Linearity, the second distinctive ingredient, is tracked only
to enforce the simulator-state-versus-operator-parameter ownership split
(\cref{app:calc:ownership}); it rides on the \code{Sim} monad's \code{get}/\code{put}
discipline and the \code{!} annotation rather than on per-tensor annotations.

\section{Evaluation order and confluence}
\label{app:calc:confluence}

A final, reassuring property. The reduction rules above can often fire in more
than one place in a term---in \cref{fig:appA:reduction}, we could have evaluated a
sub-expression before or after another. Does the order matter? \emph{No, beyond
what data dependency forces.} Because every function is pure---its output depends
only on its inputs, with no side effects (\cref{ch:language})---the value of an
expression depends only on the values of its sub-expressions, never on the order
in which those sub-expressions were reduced. Two redexes that do not depend on
each other may be fired in either order, or simultaneously, and reach the same
value; this is the calculus's \emph{confluence}. The only ordering the dynamics
respects is the one the data edges demand: a step that consumes the residual must
wait for the step that produces it.

Two qualifications keep this honest. First, demand-driven pruning
(\cref{app:calc:pruning}) means some redexes are never fired at all---an unobserved
output's sub-computation is pruned rather than reduced---but this does not disturb
the observed result, since a pruned node by definition feeds nothing observed.
Second, the calculus is deterministic only \emph{up to floating-point
reassociation}: when a load-bearing numerical trick fixes a particular reduction
or accumulation order for a property it buys (determinism, a condition number, IEEE
compliance), that order is part of the algorithm and is preserved as an explicit
algebraic claim, not freely reassociated (the load-bearing-versus-transparent
distinction of the main text). With that caveat, the program text \emph{is} the
dependency graph: there is no hidden sequencing, only data flowing along visible
edges, exactly as \cref{ch:language} promised.

\summarysection
This appendix is the formal companion to the main text's calculus chapters. The
\emph{grammar} fixes the legal syntax in two layers: a \emph{type language} $\tau$
(\cref{app:calc:types}: ground scalars, \code{Tensor[$\sigma$]}, records, tuples,
functions, \code{Sim $\tau$}, \code{Op[$\cdot$]}, and the \code{!} shareable
annotation) and a \emph{term language} $e$ (\cref{app:calc:terms}: the lambda core,
records, tuples, conditionals, primitive and operator application, the monad and
state effects, and the single loop \code{while\_loop}). The \emph{small-step
operational semantics} (\cref{app:calc:reduction}) is the dynamics, a reduction
relation $e \to e'$ presented as inference rules: $\beta$ and \code{let} for the
lambda core; projection, destructuring, and spread for records; the three monad
laws and the state-effect rewrites for orchestration; one substitution rule for
operator \code{apply}; the recursive unrolling rule for \code{iterate\_while}; and
the demand-driven pruning rule. The grammar says what you may write; the reduction
rules say what it means; and the whole computational content is the handful of
rules collected here. The signature move is \emph{pruning}: a non-observation
premise licenses replacing a producer by the subgraph of its observed outputs, so
diagnostics and production are one program. Because every function is pure, the
calculus is confluent (\cref{app:calc:confluence})---evaluation order beyond data
dependency is unobservable, deterministic up to deliberate floating-point
reassociation---so the program text is exactly the dependency graph. The ownership
split (\cref{app:calc:ownership}) and the standard typing judgment
(\cref{app:calc:typing}) round out the static side; tensor shapes are handed off
to \cref{app:shapes} and primitive $\delta$-rules and algebraic laws to
\cref{app:laws}.

\furtherreading
The calculus is a small dialect of the typed lambda calculus extended with a state
monad and one iteration combinator; the two most readable entry points are the
standard reference for the language whose grammar and reduction conventions we
borrow~\citep{peytonjones2003}, and Wadler's exposition of how a pure calculus
nonetheless models stateful computation through monads~\citep{wadler1995monads},
which underwrites both the \code{Sim} effects of \cref{app:calc:state} and the
\code{do}-notation desugaring of \cref{app:calc:monad}. Readers who want the
operational-semantics machinery itself---inference rules, $\beta$-reduction,
substitution, confluence---in its original setting will find it in any text on the
lambda calculus or programming-language theory. The two companion appendices carry
the pieces this one hands off: \cref{app:shapes} formalizes the shape grammar
$\sigma$ and its congruence algebra, and \cref{app:laws} collects the primitive
$\delta$-rules and the algebraic equational laws the reduction chains of
\cref{ch:rotations} appeal to.
